
\part{Gaussians}

\chapter{Gaussian distribution}

A continuous probability distribution, defined by two parameters: its mean \( \mu \) and standard deviation \( \sigma \).

\paragraph{PDF:}
\begin{equation*}
    f(x \mid \mu, \sigma)
    =
    \frac{1}{\sigma \sqrt{2\pi}}
    \exp \left\{ -\frac{{(x - \mu)}^2}{2\sigma^2} \right\}
\end{equation*}

\section{Properties}

\begin{itemize}
    \item The mean, median, and mode are all equal.

    \item Symmetric about the mean.
\end{itemize}

\chapter{Gaussian mixture models (GMMs)}

A probabilistic model that assumes that the data is generated from a mixture of several Gaussian distributions with unknown parameters.
A GMM is defined as a weighted sum of \( K \) Gaussian components.

\paragraph{PDF:}
\begin{equation*}
    p(x)
    =
    \sum_{k=1}^{K} \pi_k \mathcal{N}(x \mid \mu_k, \Sigma_k)
\end{equation*}

\begin{itemize}
    \item \( K \) is the number of Gaussian components.

    \item \( \pi_k \) is the mixing coefficient for the \( k \)-th component, satisfying \( \sum_{k=1}^{K} \pi_k = 1 \) and \( \pi_k \geq 0 \).

    \item \( \mathcal{N}(x \mid \mu_k, \Sigma_k) \) is the Gaussian distribution with mean \( \mu_k \) and covariance matrix \( \Sigma_k \):
        \begin{equation*}
            \mathcal{N}(x \mid \mu_k, \Sigma_k)
            =
            \frac{1}{{(2\pi)}^{d/2} |\Sigma_k|^{1/2}}
            \exp \left\{ -\frac{1}{2}{(x - \mu_k)}^\top \Sigma_k^{-1} (x - \mu_k) \right\}
        \end{equation*}

    \item \( d \) is the dimensionality of the data.
\end{itemize}

\section{Parameter estimation}

The parameters of a GMM, including the mixing coefficients \( \pi_k \), means \( \mu_k \), and covariance matrices \( \Sigma_k \), can be estimated using the \textit{Expectation-Maximization (EM)} algorithm.
The EM algorithm iteratively refines the parameter estimates to maximize the likelihood of the observed data.

\subsection{Expectation-Maximization (EM) algorithm}

The EM algorithm consists of two main steps: the Expectation (E) step and the Maximization (M) step.
\begin{itemize}
    \item \textbf{E-step}:
        Compute the posterior probabilities (responsibilities) that each data point \( x_i \) belongs to each Gaussian component \( k \):
        \begin{equation*}
            \gamma_{ik}
            =
            \frac{\pi_k \mathcal{N}(x_i \mid \mu_k, \Sigma_k)}
            {\sum_{j=1}^{K} \pi_j \mathcal{N}(x_i \mid \mu_j, \Sigma_j)}
        \end{equation*}

    \item \textbf{M-step}:
        Update the parameters using the responsibilities computed in the E-step:
        \begin{align*}
            N_k
            & =
            \sum_{i=1}^{N} \gamma_{ik}
            \\
            \pi_k^{\text{new}}
            & =
            \frac{N_k}{N}
            \\
            \mu_k^{\text{new}}
            & =
            \frac{1}{N_k} \sum_{i=1}^{N} \gamma_{ik} x_i
            \\
            \Sigma_k^{\text{new}}
            & =
            \frac{1}{N_k}
            \sum_{i=1}^{N} \gamma_{ik} (x_i - \mu_k^{\text{new}}) {(x_i - \mu_k^{\text{new}})}^\top
        \end{align*}
        where \( N \) is the total number of data points.
\end{itemize}

The E-step and M-step are repeated until convergence, typically when the change in the log-likelihood of the data given the model parameters is below a certain threshold.
